\uuid{BgeC}
\titre{Système linéaire à un paramètre/Méthode du Pivot de Gauss}
\auteur{hanani ; mzari}
\organisation{exo7}
\licence{CC BY-NC-SA 4.0}
\datecreate{2020-01-20}
\niveau{L1}
\module{Algèbre}
\chapitre{Déterminant, système linéaire}
\sousChapitre{Système linéaire, rang}
\difficulte{3}
\theme{Systèmes d'équations linéaires, Système linéaire à un paramètre, Méthode du Pivot de Gauss}
\exo7qcmid{lille-2.systemes.012}

\contenu{
\begin{question}
\qclassification{200.03}
\qtags{motcle=Système linéaire à un paramètre/Méthode du Pivot de Gauss}
On considère le système d'équations, d'inconnue $(x,y,z)\in \Rr^3$ et de paramètre un réel $m$ :
$$(\mathtt{S})\left\{\begin{array}{rcc}
x-y-z&=&1\\ -x+2y-mz&=&-3\\ 2x-y+(m-1)z&=&2m+2.\end{array}\right.$$
Quelles sont les assertions vraies ?
\begin{answers}
\good{$(\mathtt{S}) \Leftrightarrow
\left\{\begin{array}{rcc}x-y-z&=&1\\y-(m+1)z&=&-2\\
(m+1)z&=&m+1.\end{array}\right.$}
\bad{Pour tout réel $m$, $(\mathtt{S})$ admet une infinité de solutions.}
\bad{Si $m=-1$, $(\mathtt{S})$ n'admet pas de solution.}
\good{Si $m \neq -1$, $(\mathtt{S})$ admet une unique solution.}
\end{answers}
\begin{explanations}
$$(\mathtt{S}) \Leftrightarrow
\left\{\begin{array}{rcc}
x-y-z&=&1\\
y-(m+1)z&=&-2\\
y+(m+1)z&=&2m\\
\end{array}\right.  \Leftrightarrow
\left\{\begin{array}{rcc}
x-y-z&=&1\\
y-(m+1)z&=&-2\\
(m+1)z&=&m+1.\\
\end{array}\right.$$
Si $m=-1$, $(\mathtt{S})$ admet une infinité de solutions : $\{(-1+z,-2,z) \, ; \, z \in \Rr\}$.\\
Si $m \neq -1$, $(\mathtt{S})$ admet une unique solution : $(1+m,-1+m,1)$.
\end{explanations}
\end{question}
}
